\subsection*{最初に必要な用語}
\addcontentsline{toc}{subsection}{最初に必要な用語}

\noindent\par\smallskip\noindent\refstepcounter{table}\label{tab:ch12-01}表 \thetable: 最初に必要な用語の整理\par\smallskip
表 \thetable は、最初に必要な用語の整理を比較・確認しやすい形で整理したものである。\par\smallskip
\begin{tabularx}{\textwidth}{L{0.27\textwidth}Y}
\toprule
用語 & 初学者向けの説明 \\
\midrule
ソフトウェア品質 & ソフトウェアが、明示された要求と暗黙のニーズを、指定された条件のもとでどの程度満たすかを表す概念である。 \\
\term{品質要求}{Quality Requirement} & 性能、信頼性、安全性、セキュリティ、使いやすさなど、機能の実現水準や制約を表す要求である。 \\
品質マネジメント & 組織やプロジェクトが品質を方向づけ、管理するための活動である。 \\
品質保証 & プロセスや成果物が適切な品質を生み出すという信頼を得るための活動である。 \\
品質制御 & 成果物を測定・評価し、欠陥や問題を見つける活動である。 \\
検証 & 作ったものが、前工程や仕様で課された条件を満たしているかを確認することである。 \\
妥当性確認 & 作ったものが、意図した利用目的と利用者ニーズに合っているかを確認することである。 \\
\term{ディペンダビリティ}{Dependability} & 信頼して使える性質のまとまりであり、可用性、信頼性、保守性、安全性、セキュリティなどを含む。 \\
\bottomrule
\end{tabularx}
